% Options for packages loaded elsewhere
% Options for packages loaded elsewhere
\PassOptionsToPackage{unicode}{hyperref}
\PassOptionsToPackage{hyphens}{url}
\PassOptionsToPackage{dvipsnames,svgnames,x11names}{xcolor}
%
\documentclass[
  letterpaper,
  DIV=11,
  numbers=noendperiod]{scrartcl}
\usepackage{xcolor}
\usepackage{amsmath,amssymb}
\setcounter{secnumdepth}{-\maxdimen} % remove section numbering
\usepackage{iftex}
\ifPDFTeX
  \usepackage[T1]{fontenc}
  \usepackage[utf8]{inputenc}
  \usepackage{textcomp} % provide euro and other symbols
\else % if luatex or xetex
  \usepackage{unicode-math} % this also loads fontspec
  \defaultfontfeatures{Scale=MatchLowercase}
  \defaultfontfeatures[\rmfamily]{Ligatures=TeX,Scale=1}
\fi
\usepackage{lmodern}
\ifPDFTeX\else
  % xetex/luatex font selection
\fi
% Use upquote if available, for straight quotes in verbatim environments
\IfFileExists{upquote.sty}{\usepackage{upquote}}{}
\IfFileExists{microtype.sty}{% use microtype if available
  \usepackage[]{microtype}
  \UseMicrotypeSet[protrusion]{basicmath} % disable protrusion for tt fonts
}{}
\makeatletter
\@ifundefined{KOMAClassName}{% if non-KOMA class
  \IfFileExists{parskip.sty}{%
    \usepackage{parskip}
  }{% else
    \setlength{\parindent}{0pt}
    \setlength{\parskip}{6pt plus 2pt minus 1pt}}
}{% if KOMA class
  \KOMAoptions{parskip=half}}
\makeatother
% Make \paragraph and \subparagraph free-standing
\makeatletter
\ifx\paragraph\undefined\else
  \let\oldparagraph\paragraph
  \renewcommand{\paragraph}{
    \@ifstar
      \xxxParagraphStar
      \xxxParagraphNoStar
  }
  \newcommand{\xxxParagraphStar}[1]{\oldparagraph*{#1}\mbox{}}
  \newcommand{\xxxParagraphNoStar}[1]{\oldparagraph{#1}\mbox{}}
\fi
\ifx\subparagraph\undefined\else
  \let\oldsubparagraph\subparagraph
  \renewcommand{\subparagraph}{
    \@ifstar
      \xxxSubParagraphStar
      \xxxSubParagraphNoStar
  }
  \newcommand{\xxxSubParagraphStar}[1]{\oldsubparagraph*{#1}\mbox{}}
  \newcommand{\xxxSubParagraphNoStar}[1]{\oldsubparagraph{#1}\mbox{}}
\fi
\makeatother


\usepackage{longtable,booktabs,array}
\usepackage{calc} % for calculating minipage widths
% Correct order of tables after \paragraph or \subparagraph
\usepackage{etoolbox}
\makeatletter
\patchcmd\longtable{\par}{\if@noskipsec\mbox{}\fi\par}{}{}
\makeatother
% Allow footnotes in longtable head/foot
\IfFileExists{footnotehyper.sty}{\usepackage{footnotehyper}}{\usepackage{footnote}}
\makesavenoteenv{longtable}
\usepackage{graphicx}
\makeatletter
\newsavebox\pandoc@box
\newcommand*\pandocbounded[1]{% scales image to fit in text height/width
  \sbox\pandoc@box{#1}%
  \Gscale@div\@tempa{\textheight}{\dimexpr\ht\pandoc@box+\dp\pandoc@box\relax}%
  \Gscale@div\@tempb{\linewidth}{\wd\pandoc@box}%
  \ifdim\@tempb\p@<\@tempa\p@\let\@tempa\@tempb\fi% select the smaller of both
  \ifdim\@tempa\p@<\p@\scalebox{\@tempa}{\usebox\pandoc@box}%
  \else\usebox{\pandoc@box}%
  \fi%
}
% Set default figure placement to htbp
\def\fps@figure{htbp}
\makeatother





\setlength{\emergencystretch}{3em} % prevent overfull lines

\providecommand{\tightlist}{%
  \setlength{\itemsep}{0pt}\setlength{\parskip}{0pt}}



 


\KOMAoption{captions}{tableheading}
\makeatletter
\@ifpackageloaded{caption}{}{\usepackage{caption}}
\AtBeginDocument{%
\ifdefined\contentsname
  \renewcommand*\contentsname{Table of contents}
\else
  \newcommand\contentsname{Table of contents}
\fi
\ifdefined\listfigurename
  \renewcommand*\listfigurename{List of Figures}
\else
  \newcommand\listfigurename{List of Figures}
\fi
\ifdefined\listtablename
  \renewcommand*\listtablename{List of Tables}
\else
  \newcommand\listtablename{List of Tables}
\fi
\ifdefined\figurename
  \renewcommand*\figurename{Figure}
\else
  \newcommand\figurename{Figure}
\fi
\ifdefined\tablename
  \renewcommand*\tablename{Table}
\else
  \newcommand\tablename{Table}
\fi
}
\@ifpackageloaded{float}{}{\usepackage{float}}
\floatstyle{ruled}
\@ifundefined{c@chapter}{\newfloat{codelisting}{h}{lop}}{\newfloat{codelisting}{h}{lop}[chapter]}
\floatname{codelisting}{Listing}
\newcommand*\listoflistings{\listof{codelisting}{List of Listings}}
\makeatother
\makeatletter
\makeatother
\makeatletter
\@ifpackageloaded{caption}{}{\usepackage{caption}}
\@ifpackageloaded{subcaption}{}{\usepackage{subcaption}}
\makeatother
\usepackage{bookmark}
\IfFileExists{xurl.sty}{\usepackage{xurl}}{} % add URL line breaks if available
\urlstyle{same}
\hypersetup{
  pdftitle={Syllabus},
  colorlinks=true,
  linkcolor={blue},
  filecolor={Maroon},
  citecolor={Blue},
  urlcolor={Blue},
  pdfcreator={LaTeX via pandoc}}


\title{Syllabus}
\usepackage{etoolbox}
\makeatletter
\providecommand{\subtitle}[1]{% add subtitle to \maketitle
  \apptocmd{\@title}{\par {\large #1 \par}}{}{}
}
\makeatother
\subtitle{BUS 411 --- Fixed Income Security Analysis and Valuation}
\author{}
\date{}
\begin{document}
\maketitle


\section{Course information}\label{course-information}

\textbf{Instructor:} Ali Boloor\\
\textbf{Semester:} Summer 2026\\
\textbf{Email:}
\href{mailto:aboloorf@sfu.ca}{\nolinkurl{aboloorf@sfu.ca}}\\
\textbf{LMS:} \href{https://canvas.sfu.ca/}{Canvas} ,
\href{https://aliboloor.github.io/BUS411/}{Course website}\\
\textbf{Lectures:} Tuesdays, 2:30-5:20pm\\
\textbf{Office hours:} \emph{To be announced}

\section{Course Description}\label{course-description}

This course introduces the market valuation of fixed income securities.
Topics include the history of fixed income markets; valuation methods
for fixed and variable annuities, government and corporate bonds, and
bonds with contingencies; credit risk modeling; fixed income
derivatives; and risk management using duration and convexity.

\textbf{Prerequisite:} BUS 315 and BUS 360W, each with a minimum grade
of C-; 60 units.

\section{Textbook}\label{textbook}

Fabozzi, F. J. (2015). \emph{Bond markets, analysis, and strategies}
(9th ed.). Pearson. ISBN 978-0-13-379677-3.

\section{Class Format}\label{class-format}

Classes are primarily lecture-based; discussion is encouraged. Materials
include slides and lecture notes. You are expected to review the
assigned readings and textbook chapters before each class.

\section{Assessment}\label{assessment}

\begin{longtable}[]{@{}lr@{}}
\toprule\noalign{}
Component & Weight \\
\midrule\noalign{}
\endhead
\bottomrule\noalign{}
\endlastfoot
Exam 1 & 35\% \\
Exam 2 & 35\% \\
Assignment 1 & 15\% \\
Assignment 2 & 15\% \\
\end{longtable}

\section{Topics}\label{topics}

\begin{enumerate}
\def\labelenumi{\arabic{enumi}.}
\tightlist
\item
  Introduction
\item
  Bond pricing and yields
\item
  Bond price volatility
\item
  Duration, convexity, and portfolio management
\item
  Term structure and curve strategies
\item
  Interest rate models
\item
  Mortgage market
\item
  Treasury securities and corporate debt instruments
\item
  Embedded options
\item
  Corporate bond credit analysis
\item
  Credit risk modeling
\item
  Interest rate derivatives
\end{enumerate}

\section{Academic Honesty}\label{academic-honesty}

Plagiarism is the unacknowledged use of other people's ideas or work. It
is often unintentional and can be avoided through careful habits and
familiarity with academic conventions. Whether intentional or not,
plagiarism is a serious academic offence. The university's stance
reflects a shared commitment to intellectual honesty; original work by
students and faculty sustains the university as a centre of knowledge
and research. It is your responsibility to acknowledge and cite all
resources you use.

Representative examples of academic dishonesty include:

\begin{itemize}
\tightlist
\item
  Presenting another person's work as your own (plagiarism).
\item
  Submitting the same work more than once without prior approval.
\item
  Translating a work from one language to another without complete and
  proper citation.
\item
  Cheating.
\item
  Impersonation (e.g., having someone else write your exam).
\item
  Submitting false records or information (e.g., forged medical notes).
\item
  Stealing or destroying another student's work.
\item
  Unauthorized or inappropriate use of computers, phones, calculators,
  or other technology in coursework, assignments, or examinations.
\item
  Falsifying material subject to academic evaluation.
\item
  Any activity intended to circumvent standards of academic honesty,
  even if not listed here.
\end{itemize}

You are expected to post comments and to write reports and exams in your
own words. Whenever you use an idea or wording from another source, cite
it appropriately. If you are struggling with an assignment, contact the
instructor or the program office for help.

Ignorance of these standards does not excuse academic dishonesty.

The full SFU policy on academic honesty is available here:
\href{https://www.sfu.ca/policies/gazette/student.html}{SFU policies ---
student academic integrity}.




\end{document}
